% =============================================================================
% main.tex — Thin-root document for the agentic-publishing-pipeline report.
%
% Declares the document class (report, for \chapter-level structure) and
% \input{}s every other file.  No narrative text, table environments,
% TikZ pictures, or inline equations appear here.
%
% Compiler: LuaLaTeX + biber + makeindex (two-pass build)
%   lualatex main.tex
%   biber main
%   makeindex main.nlo -s nomencl.ist -o main.nls
%   makeindex main
%   lualatex main.tex
%   lualatex main.tex
% =============================================================================

\documentclass[11pt,a4paper]{report}

% --- File-existence guards ---------------------------------------------------
% Defined immediately so they are available before the preamble/macros inputs.
%
% \RequireAndInput{stem}   — inputs stem.tex from the document root;
%                            errors clearly if the file is missing.
% \RequireAndChapter{stem} — inputs chapters/stem.tex and increments a
%                            counter so a deleted line is caught below.
%
% Bug-fix notes vs. original:
%  * Definitions now precede the first \RequireAndInput call (ordering fix).
%  * \IfFileExists check uses the .tex extension explicitly for portability.
%  * \input is inside the true branch only — prevents the double-error that
%    occurs in nonstop mode when both \PackageError and \input fire for the
%    same missing file.
%  * Counter approach replaces \csname chapter#1true\endcsname so chapter
%    stems containing underscores (e.g. tool_use) work without issue.
\newcommand{\RequireAndInput}[1]{%
  \IfFileExists{#1.tex}{%
    \input{#1}%
  }{%
    \PackageError{main}{Required file '#1.tex' not found}{%
      The file '#1.tex' is required by this document but could not be
      located.  Check that the file exists and the path is correct.%
    }%
  }%
}

\newcounter{chaptersexpected}
\newcounter{chaptersincluded}
\setcounter{chaptersexpected}{8}% must equal the number of \RequireAndChapter calls

\newcommand{\RequireAndChapter}[1]{%
  \IfFileExists{chapters/#1.tex}{%
    \stepcounter{chaptersincluded}%
    \input{chapters/#1.tex}%
  }{%
    \PackageError{main}{Required chapter '#1' not found}{%
      'chapters/#1.tex' is required by this document but could not be
      located.%
    }%
  }%
}

% --- Preamble and shared macros ----------------------------------------------
% These must come after the \newcommand definitions above.
\RequireAndInput{preamble}
\RequireAndInput{macros}

% --- Title-page metadata -----------------------------------------------------
\title{Agentic Reasoning for Large Language Models: A Survey}
\author{Team Name \\ Course: AI Orchestration}
\date{\today}

% =============================================================================
\begin{document}
% =============================================================================

% --- Front matter ------------------------------------------------------------
\maketitle
\tableofcontents

% --- Main body (chapters) ----------------------------------------------------
% Chapter order matches constants.py::REQUIRED_CHAPTER_IDS.
% A missing file produces a \PackageError immediately; a deleted
% \RequireAndChapter line is caught by the counter guard below.
\RequireAndChapter{introduction}
\RequireAndChapter{planning}
\RequireAndChapter{memory}
\RequireAndChapter{retrieval}
\RequireAndChapter{tool_use}
\RequireAndChapter{multimodal}
\RequireAndChapter{evaluation}
\RequireAndChapter{conclusion}

% --- Chapter-count guard (runs BEFORE back matter) ---------------------------
% Fires before \printnomenclature etc. so that in nonstop mode a deleted
% \RequireAndChapter line prevents back-matter output as well.
\ifnum\value{chaptersincluded}=\value{chaptersexpected}\else
  \PackageError{main}{%
    Chapter count mismatch: \the\value{chaptersincluded}\space of
    \the\value{chaptersexpected}\space required chapters were included.
    A \string\RequireAndChapter{} line may have been deleted.%
  }{}%
\fi

% --- Back matter -------------------------------------------------------------
% TOC entries for nomenclature, bibliography, and index are added
% automatically by the intoc options in preamble.tex.  No \addcontentsline
% is needed here.
\printnomenclature
\printbibliography[heading=bibintoc]
\printindex

% =============================================================================
\end{document}
